\label{sec:approach}
The project was approached as two interacting systems: one interprets photographed handwriting and constructs writer-specific paths, while the other converts those paths into physical pen movement. The proposal's functional analysis provides their responsibilities; the G-code interface allows their implementation and verification to proceed separately.

\subsection{Mechanical Hardware Design}

Before any computer vision or neural networks were worked on, the hardware to perform the pen plotting was created. This is because the hardware does not rely on any of the coding elements to be complete to be able to be tested, and hardware can experience major delays in production due to order delays, stock shortages or component failure, adding to the reason for it to be constructed first. A Cartesian gantry with two independent axes was chosen over some sort of robotic writing arm, as system requiring only flat, cartesian movement does not need joint-dependent kinematics, and each axis can be calibrated separately and work indepedently from one another. To achieve the linear motion required, GT2 timing belts and pulleys were chosen over a lead-screws due to their lower cost, faster travel and ease of customizstion. 

Two NEMA 17 stepper motors were chosen for the X and Y axes, driven by A4988 driver chips. Because a stepper motor always rotates by precisely the number of commanded steps at exact intervals, accuracy could be assured without an encoder. To allow for even more increased accuracy and minimal friction, linear bearings were used within the mechanics of the system to allow parts to move with ease and precision, instead of just having the printed supports sliding over the steel rails of the gantry. The use of microswitches on the axes allowed the system to be able to calibrate to reproduce text more accurately. The gantry was designed to be able to reproduce any G-code file that was provided, allowing for easy integration with the other subsystems in future.

\subsection{Page condition and line segmentation}

In parallel while working on the gantry and waiting for hardware to be delivered, work was done on the first processing stage of the software pipeline, which was performing image conditioning and segmentation of those pages into a format the future classifiers would accept. A photograph produces a wide range of inconsistencies that need to be dealt with, including uneven lighting, page rotation, printed rules, and strokes from neighbouring lines merging into each other. The image conditioning part of the code was designed to remove these undesired page inconsistencies without ruining the quality of the original writing, using methods like lighting normalization, ink separation, and skew correction. 

Segmentation of the page of text was chosen to be done lines at a time instead of per word or character, as it was found that dynamic word separation that worked in all cases was difficult, and line separation has the added advantage of more context to train and classify text. The design chosen to perform the segmentation involves locating line centres first, finds curved boundaries between the lines and its ascenders and descenders, and only then uses the geometry of each ink piece to decide which line actually owns it. Deskewing, removal of image background, and artefact rejection were deliberately ordered to run before line centres are estimated, while recovery of faint ink was left until after ownership of each piece had already been decided, since doing it earlier risked reattaching noise to the wrong line. The system fomatted each page into the ideal crop sizes to be passed into the text and author classifiers.

\subsection{Text and writer classification}

Once clean line crops were available, the next stage had two jobs: reading the text and identifying who wrote it. Both depend on the same stroke features, so a single convolutional backbone was built, but then the writer and author classifiers varied after that point. The text classifier uses a CNN–BiLSTM architecture with decoded CTC output, which was heavily based off existing research, and chosen because it reads a whole line without needing individual characters pre-cut, as mentioned in the segmentation section, and cursive text has no reliable character boundary to cut at in the first place.

Training data came from a mix of public handwriting (IAM) and the project’s own photographed pages, and both author and text models were trained using hold-out pages or lines to allow the system to be validated against, to prove over-fitting did not occur. The writer network’s input was stroke-normalised rather than using varying ink thickness from real pages, specifically because the gantry’s own pen produces constant stroke thicknesses, regardless of the author.

\subsection{Style extraction and handwriting reproduction}

 Once the previous subsystems were created, the actual text synthesis and reproduction was made. Four reproduction methods were considered, and the method chosen instead was a stored set of normalised glyph trajectories. The other were not chosen because an online dataset of pen strokes was not available, and the chosen IAM dataset was not extensive enough to train a rpeproduction model, neither was the necessary GPU processing power available for it. Using the glyph trajectory method simply resused data and models the project already had, but at the cost of weaker handling of how letters join and vary with context compared to a true generative model especially for heavy cursive styles. The new text to be created is then assembled from the stored glyphs rather than generated from scratch, giving direct, inspectable control over scale and spacing before anything is sent to the gantry. Extracted features like slant, spacing etc. are also applied to the synthesised text.

\subsection{Integration and testing}

The integration stage formed a complete pipeline of the system through a web server, with G-code remaining the constant fixed interface between the software reproduction output and the gantry hardware. The web interface was designed to be able to preview and inspect generated G-cocdes before any gantry operation was called. By previewing the gantry before reproduction, the precision of the gantry could be measured, and any errors with gantry movement isolated. Safety features were also included as part of the integration, which included a Emergency abort feature on the web server to manually stop any gantry command, as well as end-stop detection with switches to ensure the motors did not stall if the end of the rail was ever reached. 